\begin{example}
\label{example-idempotent-matrices}
For another example, consider
$R = k[\{t_{ij}\}_{i, j = 1}^{n}]/\mathfrak a$, where $\mathfrak a$ is
the ideal generated by the entries of the product matrix $T^2-T$,
$T = (t_{ij})$. From linear algebra, we know that under the
$GL(n, k)$-action defined by $g, T \mapsto gTg^{-1}$, $T$ is
classified by the its rank and each $T$ is conjugate to some
$\text{diag}(1, \ldots, 1, 0, \ldots, 0)$, which has $r$ 1's and $n-r$ 0's.
Thus each orbit of such a $\text{diag}(1, \ldots, 1, 0, \ldots, 0)$ under the
group action forms an irreducible component and every idempotent
matrix is contained in one such orbit. Next we will show that any
two different orbits are necessarily disjoint. For this purpose we
only need to cook up polynomial functions that take different
values on different orbits. In characteristic 0 cases, such a
function can be taken to be
$f(t_{ij}) = trace(T) = \sum_{i = 1}^nt_{ii}$. In positive
characteristic cases, things are slightly more tricky since we
might have $trace(T) = 0$ even if $T \neq 0$. For instance, $char = 3$
$$
trace\left(
\begin{matrix}
1 & & \\
& 1 & \\
& & 1
\end{matrix}
\right) = 3 = 0
$$
Anyway, these components can be separated using other functions.
For instance, in the characteristic 3 case, $tr(\wedge^3T)$ takes
value 1 on the components corresponding to $diag(1, 1, 1)$ and 0 on
other components.
\end{example}